\section{SONUÇLAR VE ÖNERİLER}
\label{bolum10}

Bu proje kapsamında geliştirilen API Fortress sistemi, REST API güvenliği alanında uygulamalı eğitim sağlayan saldırı ve savunma simülasyon laboratuvarı olarak başarıyla tamamlanmıştır. Sistem içerisinde hem güvenlik açıkları içeren Vulnerable API ortamı hem de modern savunma mekanizmalarıyla güçlendirilmiş Secure API ortamı birlikte geliştirilmiştir. Böylece kullanıcıların yalnızca teorik güvenlik bilgisi edinmesi değil; aynı zamanda gerçek saldırı senaryolarını uygulamalı olarak deneyimleyebilmesi ve güvenli backend mimarisinin nasıl çalıştığını doğrudan inceleyebilmesi sağlanmıştır.

Proje kapsamında OWASP API Security Top 10 içerisinde yer alan modern API güvenlik problemleri temel alınmıştır. Özellikle Mass Assignment, Broken Authentication, IDOR/BOLA ve Broken Function Level Authorization gibi kritik güvenlik açıkları laboratuvar ortamına başarılı şekilde entegre edilmiştir. Bu sayede kullanıcılar yetkisiz erişim, role escalation, authentication bypass ve object-level authorization problemleri gibi gerçek dünya saldırı senaryolarını kontrollü bir ortamda analiz edebilmiştir.

Güvenli API ortamında ise modern REST API güvenlik prensiplerini temsil eden çeşitli savunma mekanizmaları geliştirilmiştir. JWT tabanlı kimlik doğrulama sistemi ile kullanıcı doğrulama süreçleri güvenli hale getirilmiş, rol tabanlı yetkilendirme (RBAC) sistemi ile kullanıcı ve yönetici işlemleri birbirinden ayrılmıştır. Ayrıca Flask-Limiter kullanılarak brute force saldırılarına karşı rate limiting mekanizması uygulanmış, regex tabanlı WAF sistemi ile SQL Injection ve XSS payload’larının filtrelenmesi sağlanmıştır. Kullanıcı parolalarının bcrypt algoritmasıyla hashlenmesi ve input validation süreçlerinin uygulanması sayesinde veri güvenliği artırılmıştır.

Gerçekleştirilen testler sonucunda Vulnerable API ortamında oluşturulan zafiyetlerin başarıyla sömürülebildiği doğrulanmıştır. Secure API ortamında ise geliştirilen savunma mekanizmalarının büyük ölçüde etkili olduğu gözlemlenmiştir. Özellikle JWT doğrulama sistemi, authorization kontrolleri, request filtering yapıları ve rate limiting mekanizmalarının sistem güvenliğini önemli ölçüde artırdığı görülmüştür. Böylece proje, saldırı ve savunma mekanizmalarının aynı sistem üzerinde karşılaştırmalı biçimde incelenebildiği uygulamalı bir laboratuvar ortamı sunmuştur.

Bu çalışma ile birlikte elde edilen en önemli katkılardan biri, REST API güvenliği üzerine doğrudan odaklanan özgün bir eğitim platformunun geliştirilmiş olmasıdır. Mevcut birçok güvenlik laboratuvarı genel web güvenliğine yoğunlaşırken API Fortress özellikle REST API zafiyetlerini hedef almış ve bu açıkları saldırı ve savunma boyutlarıyla birlikte sunmuştur. Böylece kullanıcılar modern backend sistemlerinde karşılaşılan güvenlik problemlerini daha gerçekçi biçimde analiz edebilmiştir.

Projenin önemli katkılarından biri de Docker tabanlı taşınabilir laboratuvar mimarisidir. Konteyner tabanlı yapı sayesinde sistem farklı işletim sistemlerinde ve geliştirme ortamlarında kolaylıkla çalıştırılabilmektedir. Bu yaklaşım hem laboratuvar ortamının tekrar üretilebilirliğini artırmış hem de eğitim erişilebilirliğini kolaylaştırmıştır. Kullanıcılar karmaşık kurulum süreçleriyle uğraşmadan sistemi çalıştırarak güvenlik testlerini gerçekleştirebilmektedir.

CTF (Capture The Flag) benzeri saldırı senaryolarının sisteme entegre edilmesi de aktif öğrenme deneyimini destekleyen önemli bir unsur olmuştur. Kullanıcılar yalnızca teorik bilgi okumak yerine gerçek güvenlik açıklarını çözmeye çalışarak sistem davranışını analiz edebilmiştir. Bu yaklaşım, güvenlik farkındalığının artırılması ve öğrenme sürecinin daha kalıcı hale gelmesi açısından önemli avantaj sağlamıştır.

Ayrıca güvenli ve güvensiz API ortamlarının aynı sistem içerisinde birlikte sunulması sayesinde karşılaştırmalı analiz imkânı oluşturulmuştur. Kullanıcılar aynı endpoint’in güvensiz sürümünde saldırının neden başarılı olduğunu, güvenli sürümünde ise hangi savunma mekanizmaları nedeniyle engellendiğini doğrudan inceleyebilmiştir. Bu yaklaşım modern güvenli yazılım geliştirme prensiplerinin daha anlaşılır hale gelmesini sağlamıştır.

Proje sürecinde bazı geliştirme alanları da tespit edilmiştir. Özellikle regex tabanlı WAF yapısının yalnızca standart payload’ları filtreleyebilmesi, gelişmiş encoded veya obfuscated saldırılar karşısında ek güvenlik katmanlarına ihtiyaç duyulduğunu göstermiştir. Ayrıca sistem şu anda temel REST API güvenlik senaryolarına odaklanmaktadır. Gelecek çalışmalarda:

\begin{itemize}[leftmargin=*]
    \item gelişmiş WAF mekanizmaları,
    \item davranış tabanlı saldırı analizi,
    \item yapay zekâ destekli anomali tespiti,
    \item API Gateway entegrasyonu,
    \item merkezi log yönetimi,
    \item SIEM entegrasyonu,
    \item gerçek zamanlı saldırı izleme sistemleri
\end{itemize}

gibi daha ileri güvenlik bileşenlerinin sisteme eklenmesi planlanmaktadır.

Sonuç olarak API Fortress projesi, REST API güvenliğini uygulamalı olarak öğretmeyi amaçlayan kapsamlı bir saldırı ve savunma laboratuvarı olarak başarılı şekilde geliştirilmiştir. Proje hem modern REST API mimarisini hem de güncel güvenlik mekanizmalarını uygulamalı biçimde göstermekte; güvenli yazılım geliştirme farkındalığının artırılmasına katkı sağlamaktadır. Bu yönüyle sistem, hem eğitimsel hem de teknik açıdan REST API güvenliği alanında kullanılabilecek taşınabilir ve genişletilebilir bir laboratuvar platformu sunmaktadır.

\subsection{Gelecek Çalışmalar}
\label{gelecek-calismalar}

Sistemin genişletilebilirliği gözetilerek aşağıdaki iyileştirmeler planlanmaktadır:

\begin{itemize}[leftmargin=*]
    \item Yeni OWASP API Security senaryolarının (API6-API10) laboratuvara eklenmesi.
    \item Gerçek zamanlı saldırı analizi ve görselleştirme modülü geliştirilmesi.
    \item Farklı zorluk seviyelerine sahip CTF senaryolarının oluşturulması ve otomatik puanlama sistemi.
    \item Kullanıcı ilerleme takip sistemi ve öğrenme yolu yönetimi.
    \item Makine öğrenmesi tabanlı anomali tespiti ve davranışsal WAF modülü.
    \item Frontend kullanıcı arayüzünün geliştirilmesi ve görsel dashboard entegrasyonu.
\end{itemize}

